\documentclass[11pt,a4paper]{article}

\usepackage{fontspec}
\usepackage[english]{babel}
\babelfont{rm}{Times New Roman}
\usepackage{amsmath,amssymb}
\usepackage{booktabs}
\usepackage[margin=2.4cm]{geometry}
\usepackage[numbers,sort&compress]{natbib}

\title{Affine residualization for HHL:\\
methods, measured answers, and large-sample tests}
\author{L\^e V\~u Quang Minh}
\date{14 October 2026}

\begin{document}
\maketitle

\begin{abstract}
This note records the method and the measured answers for a budgeted hybrid solver of \(Ax=b\).
A classical step builds an explanation \(z\) on \(k\) eigenmodes.
An HHL circuit is then tuned to the remaining spectrum, and the pieces are recombined as \(x=z+\lVert b_{\mathrm{res}}\rVert\, y\).
On exact statevector simulation, with no device noise, solution fidelity stays at or above \(0.999\).
Standard HHL has relative error about \(0.005\) to \(0.06\), mostly from phase-estimation discretization.
When \(z\) is exact and the budget actually lowers \(\kappa_{\mathrm{eff}}\), the hybrid error falls to rounding level.
The energy ranking does not guarantee that drop: one stored system with \(N=4\), \(\kappa=1000\), and \(k=1\) still has \(\kappa_{\mathrm{eff}}\) about \(1000\).
On two million random Hermitian systems, an exact orthogonal split never exceeds \(\sqrt{2}\) by more than \(10^{-6}\), and a non-orthogonal shift exceeds it on every draw.
A learned \(z\) does not replace the eigendecomposition.
Against the Dalzell shortcut that costs about \(\kappa\ln(2\sqrt{2}/\varepsilon)\) block-encoding queries once the solution norm is known, exact deflation is cheaper only on spectra with a tight cluster of very small eigenvalues.
The note claims neither a quantum speedup nor a wall-clock advantage over a dense solve.
\end{abstract}

\section{Method}

\subsection{What is solved}
The target is a linear system \(Ax=b\).
The implementation takes \(A\) Hermitian, \(A=U\operatorname{diag}(\lambda)U^{\mathsf H}\), with coefficients \(\beta=U^{\mathsf H}b\).
A non-Hermitian matrix is embedded as
\[
\begin{bmatrix}0&A\\ A^{\mathsf T}&0\end{bmatrix},
\]
which has spectrum \(\pm\sigma_i\), and the clock register is signed.
Dimensions that are not powers of two are padded.
The code does not build a variational quantum linear solver and does not measure its gradient.
Barren plateaus are a motive for caring about \(\kappa\), not a measured quantity \citep{bravoprieto2023}.

HHL prepares a state proportional to \(A^{-1}b\) when \(A\) is sparse and \(\kappa\) is moderate \citep{harrow2009}.
The circuit used here is an exact statevector simulation.
The unitary \(e^{iAt}\) is applied as an exact \texttt{UnitaryGate}, not as a Trotter formula or a linear combination of unitaries \citep{gilyen2019}.
Compiled CX counts are therefore a trend under that expansion, not a hardware gate count.
Scaling in \(N\) for a genuine HHL advantage also needs a sparse block-encoding of \(A\) and an efficient preparation of \(\lvert b\rangle\).
Those assumptions are absent from the dense circuits in this study.
Later solvers reach \(O(\kappa\log(1/\varepsilon))\) queries \citep{costa2022,dalzell2024}.

\subsection{Affine split}
Preconditioning replaces \(A\) by \(MA\) \citep{clader2013}.
Affine residualization does not.
It splits the solution and leaves \(A\) unchanged:
\begin{equation}
  x=z+y.
  \label{eq:split}
\end{equation}
The classical vector \(z\) carries the modes assigned to the budget.
The circuit solves the rest.
With the Hermitian eigenbasis,
\begin{equation}
  z=\sum_{i\in R}\frac{\beta_i}{\lambda_i}u_i,
  \qquad
  b_{\mathrm{res}}=P_{\mathrm{keep}}(b-Az),
  \qquad
  x=z+\lVert b_{\mathrm{res}}\rVert\, y,
  \label{eq:hybrid}
\end{equation}
where \(y\) is the HHL state for the normalized residual \(b_{\mathrm{res}}/\lVert b_{\mathrm{res}}\rVert\), and \(R\) is the set of removed modes.
Equation~\eqref{eq:split} is~\eqref{eq:hybrid} after the residual norm has been absorbed into \(y\).
The projector \(P_{\mathrm{keep}}\) is required.
A leftover piece \(\eta u_j\) on a removed mode is returned by an exact solve as \((\eta/\lambda_j)u_j\).

The budget \(k\) is spent in one of two ways.
The energy score of a mode is
\begin{equation}
  E_j=\Bigl\lvert\frac{\beta_j}{\lambda_j}\Bigr\rvert.
\end{equation}
The \texttt{energy} rule puts the \(k\) largest scores into \(z\).
A small eigenvalue with a tiny overlap \(\beta_j\) need not be selected.
The \texttt{smallest} rule removes the \(k\) eigenvalues of smallest magnitude.
Only that rule forces
\begin{equation}
  \kappa_{\mathrm{eff}}
  =\frac{\max_{i\notin R}\lvert\lambda_i\rvert}{\min_{i\notin R}\lvert\lambda_i\rvert}
  =\frac{\lvert\lambda\rvert_{\max}}{\lvert\lambda\rvert_{k+1}}
  \le\kappa,
\end{equation}
when the eigenvalues are ordered by increasing magnitude.
Equality holds when \(\lvert\lambda_{k+1}\rvert=\lvert\lambda_{\min}\rvert\).
The inequality does not force \(\kappa_{\mathrm{eff}}\) near \(1\).

If \(z\) is exactly the piece of \(x\) on the removed modes, those modes are orthogonal to the remainder and
\begin{equation}
  \lVert x\rVert^2=\lVert z\rVert^2+\lVert x-z\rVert^2.
  \label{eq:pythagoras}
\end{equation}
Cauchy--Schwarz then gives
\begin{equation}
  \Lambda_z=\frac{\lVert z\rVert+\lVert x-z\rVert}{\lVert x\rVert}\le\sqrt{2},
\end{equation}
with equality when the two norms are equal.
This ceiling uses \(z\perp(x-z)\).
It is not a bound for an arbitrary explanation \(\hat z=z+\delta\).

\subsection{Two circuits when \(z\) is inexact}
Let \(\hat z=z+\delta\), and assume the residual solver is exact.
\begin{itemize}
  \item Reduced mode sets \(b_{\mathrm{res}}=P_{\mathrm{keep}}(b-A\hat z)\) and tunes phase estimation to \(\kappa_{\mathrm{eff}}\).
  The assembled error is \(P_{\mathrm{hard}}\delta\), so the relative error is \(\lVert P_{\mathrm{hard}}\delta\rVert/\lVert x\rVert\).
  A component of \(\delta\) in a removed mode is invisible to that circuit.
  \item Full mode keeps the raw residual \(b-A\hat z\).
  Exact arithmetic returns \(x\).
  The residual still carries the eigenvalues in \(R\), so the circuit is sized by the original \(\kappa\).
\end{itemize}
A \(10\%\) error on a removed mode is about a \(10\%\) solution error when \(\lVert z\rVert\) is close to \(\lVert x\rVert\).
The factor in general is \(\lVert z\rVert/\lVert x\rVert\).
With phase estimation, the full circuit only reaches the HHL discretization floor, about \(10^{-3}\) to \(10^{-2}\) in the trials below, and the CX count returns to the full circuit.
One illustration used \(1{,}435\) CX in the full circuit and \(185\) CX in the reduced circuit.

Clock length, evolution time, and the rotation constant \(C\) are taken from the targeted spectrum.
The clock is capped at \(12\) qubits, so \(\kappa=10^{4}\) on the small grid is resolution-limited.
Under perfect phase estimation, \(\lVert b\rVert=1\), and \(C\) no larger than the smallest targeted eigenvalue,
\begin{equation}
  P(\text{ancilla}=1)=\sum_i\lvert\beta_i\rvert^2(C/\lambda_i)^2,
\end{equation}
summed over the targeted spectrum.
Removing a mode that already sat at a full rotation can lower this probability.

\subsection{What the large tests check}
The scale tests call \texttt{bench/theory\_scale.py}.
Each draw is a Hermitian matrix with a log-uniform spectrum.
Its condition number is drawn log-uniformly from \(10\) to \(10^{4}\).

The residual is solved by \texttt{numpy.linalg.solve}.
That call stands in for an exact HHL step.
No quantum circuit is run in that test.
One random stream is used for each pair of size, budget, and ranking.
A smaller sample count is a prefix of the larger count on the same stream.

The checks on every draw are the recombination error, orthogonality of an exact split, and the ratio \(\Lambda_z\).
They also check the smallest-magnitude formula for \(\kappa_{\mathrm{eff}}\), the reduced-mode identity for a \(10\%\) hard-mode error, recovery by the raw residual, and a shift of size \(40\lVert x\rVert\) along a kept mode.
A draw is counted as stuck when \(\kappa_{\mathrm{eff}}>0.99\kappa\).

\section{Measured answers}

\subsection{Statevector grid}
Every HHL circuit below is an exact statevector, with no device noise, and \(N\le 8\) except in the query model.
The resource grid draws two random systems for each pair \((N,\kappa)\), uses the energy rule, and reports the mean.
Spectra are log-uniform on \([1/\kappa,1]\).

Fidelity \(\lvert\langle\hat x\vert x\rangle\rvert^2\) stays at or above \(0.999\) on every averaged cell.
Standard HHL (\(k=0\)) has relative error from about \(5\cdot 10^{-3}\) to \(6\cdot 10^{-2}\).
Most of that error is discretization of the solution norm, not a wrong direction of \(\lvert x\rangle\).

When \(z\) is exact and the budget brings \(\kappa_{\mathrm{eff}}\) near \(1\), the relative error falls to \(10^{-16}\)--\(10^{-14}\).
At \(N=4\), \(\kappa=1000\), standard HHL has error \(5.2\cdot 10^{-3}\); with \(k=3\) the mean hybrid error is \(8.6\cdot 10^{-14}\).
The same cell with \(k=1\) still has error \(5.3\cdot 10^{-3}\), level with standard HHL, because the energy rule did not remove the eigenvalue that sets \(\kappa\).

An emulated QSVT polynomial accurate to \(\varepsilon=10^{-2}\), on the four timed systems below, has relative error from \(6\cdot 10^{-5}\) to \(4\cdot 10^{-3}\) and fidelity from \(0.99999\).
That error is the polynomial error, not phase estimation.

At \(N=4\), \(\kappa=1000\), standard HHL uses \(12\) clock qubits, depth \(8{,}800\), and \(4{,}637\) CX gates after compilation to \(\{u,\mathrm{CX}\}\) at optimization level \(1\).
With \(k=2\), the means are about \(4.5\) clock qubits, depth \(266\), and \(159\) CX.
At \(N=8\), \(\kappa=1000\), the CX count falls from \(5{,}632\) (\(k=0\)) to \(1{,}277\) (\(k=3\)).
Absolute CX counts depend on the expansion of a general unitary.
The decrease with \(k\) is the result.

Success probability does not always rise.
On the grid at \(N=4\), \(\kappa=10\), \(k=1\), the mean probability is about \(0.18\), against about \(0.41\) at \(k=0\).

\subsection{Neural estimate of \(z\)}
A full Hermitian eigendecomposition already yields \(x_i=\beta_i/\lambda_i\).
The quantum step adds no information once those pairs are known.
An MLP was trained to predict \(z\) without that decomposition.
The median relative error \(\lVert\hat z-z\rVert/\lVert z\rVert\) is \(0.68\) at \(N=2\) (\(8{,}000\) training samples, \(1{,}000\) test samples) and \(0.95\) at \(N=4\) with the same sample size.
At \(N=8\), with \(400\) training samples, the median is \(0.99\).
Reduced-mode solution error follows this error: median \(0.67\) at \(N=2\) and \(0.84\) at \(N=4\).
There is no evidence here for an \(O(N)\) or \(O(N^2)\) predictor with a small enough error.

\subsection{Simulator wall time}
Times below are one system per cell, a fixed seed, the \texttt{smallest} rule, and no circuit compilation, in milliseconds.
\texttt{numpy.linalg.solve} stays under \(0.01\,\mathrm{ms}\) with relative error \(0\) at these sizes.

\begin{center}
\begin{tabular}{@{}llrrrrr@{}}
\toprule
\(N\) & \(\kappa\) & NumPy & HHL & Hybrid \(k=1\) & Hybrid + QSVT & Full QSVT \\
\midrule
4 & 100 & 0.007 & 266 & 53 & 0.36 & 0.43 \\
4 & 1000 & 0.005 & 1{,}834 & 89 & 0.35 & 0.67 \\
8 & 100 & 0.006 & 897 & 458 & 0.40 & 0.35 \\
8 & 1000 & 0.005 & 9{,}352 & 11{,}045 & 0.54 & 1.02 \\
\bottomrule
\end{tabular}
\end{center}

When one removed mode shortens the clock, the hybrid simulation is faster than standard HHL.
At \(N=4\), \(\kappa=1000\), the time falls from \(1.8\,\mathrm{s}\) to \(0.09\,\mathrm{s}\), and the relative error falls from \(5.2\cdot 10^{-3}\) to \(1.7\cdot 10^{-3}\).
When one mode is not enough, the simulation is not faster and can be less accurate.
At \(N=8\), \(\kappa=1000\), the hybrid run takes \(11.0\,\mathrm{s}\) against \(9.4\,\mathrm{s}\) for standard HHL, with relative error \(0.062\) against \(0.0076\) and fidelity \(0.998\).
The QSVT times are an eigenvalue polynomial, not a quantum circuit.

\subsection{Query model}
The unit is one call to a block-encoding of \(A/\lVert A\rVert\).
The Dalzell shortcut, once the solution norm is known, costs \(\kappa\ln(2\sqrt{2}/\varepsilon)\) \citep{dalzell2024}.
The comparison uses medians of \(24\) systems at \(N=64\) and \(\varepsilon=10^{-2}\).
On the log-uniform family, removing up to eight modes exactly and then running QSVT remains more expensive than that shortcut: at \(\kappa=10^{4}\), \(6.4\cdot 10^{5}\) queries against \(5.6\cdot 10^{4}\).
Textbook HHL in the same cell is \(2.4\cdot 10^{8}\).
On a spectrum with eight eigenvalues of size \(1/\kappa\) and the rest in \([0.1,1]\), removing those modes exactly leaves about \(780\) queries, against \(5.6\cdot 10^{4}\) for the shortcut.
At \(\kappa=100\) on that same clustered spectrum, the shortcut is again cheaper (\(564\) against about \(750\)).
The classical cost of the exact eigendecomposition is not included.
Once those eigenpairs are known, \(Ax=b\) is already solved.

Adaptive Lanczos, with \(k\le 8\) and at most half of \(N\) matrix--vector products, did not converge any small mode in the families tried, and the solver fell back to QSVT on the full spectrum.
Related uses of the words ``residual'' and ``affine'' denote different objects: a joint block-encoding of \(A\) and \(b\) \citep{li2025,dalzell2026}, and a Lindbladian encoding of a residual \citep{shang2026}.

\section{Large-sample tests}
\label{sec:scale}

\subsection{One million systems}
On 9 October 2026, six streams were run to \(10^{6}\) draws: \(N=2,4,8\), \(k=1\), both rankings.
Checkpoints at \(10^{4}\) and \(10^{5}\) are prefixes of the same streams.
The full table is \texttt{bench/theory\_scale\_log.csv}.

Across all one million systems, the recombination error stays at most \(5.2\cdot 10^{-12}\).
No exact split exceeds \(\sqrt{2}\) by more than \(10^{-6}\).
The smallest-magnitude formula for \(\kappa_{\mathrm{eff}}\) matches on every such system.
A \(10\%\) hard-mode error matches the reduced-mode identity to \(5.2\cdot 10^{-12}\), and the raw residual clears it below \(1.2\cdot 10^{-12}\).
A shift along a kept mode exceeds \(\sqrt{2}\) on every system.

The stuck counts, \(\kappa_{\mathrm{eff}}>0.99\kappa\), barely move as the prefix grows.

\begin{center}
\begin{tabular}{@{}llrrr@{}}
\toprule
\(N\) & rule & \(10^{4}\) & \(10^{5}\) & \(10^{6}\) \\
\midrule
2 & smallest & 0 & 0 & 0 \\
2 & energy & 0 & 0 & 0 \\
4 & smallest & 40 & 449 & 4{,}056 \\
4 & energy & 1{,}883 & 18{,}961 & 190{,}957 \\
8 & smallest & 99 & 1{,}168 & 11{,}808 \\
8 & energy & 3{,}625 & 37{,}134 & 369{,}657 \\
\bottomrule
\end{tabular}
\end{center}

At one million draws the energy rule is still stuck on \(19.1\%\) of the \(N=4\) systems and \(37.0\%\) of the \(N=8\) systems.
The smallest-magnitude rule is stuck on \(0.41\%\) and \(1.18\%\).

\subsection{Two million systems, larger than the first test}
On 13 October 2026 the same checker was run to \(2\times 10^{6}\) draws, with \(N\in\{2,4,8,16\}\) and budgets \(k=1,2\).
That is fourteen streams.
The one-million row of each stream is a prefix of the two-million row.
For \(N\in\{2,4,8\}\) and \(k=1\), that prefix matches the section above.
The raw rows are \texttt{bench/theory\_scale\_log\_2e6.csv}.

On all \(2\times 10^{6}\) systems in every stream:
\begin{itemize}
  \item recombination error stays at most \(5.9\times 10^{-12}\);
  \item no exact split has \(\Lambda_z>\sqrt{2}+10^{-6}\), and the largest stored value sits on \(\sqrt{2}\) within a unit in the last place;
  \item the smallest-magnitude formula matches on every such system, and \(\kappa_{\mathrm{eff}}\) never exceeds \(\kappa\);
  \item a \(10\%\) hard-mode error agrees with \(\lVert P_{\mathrm{hard}}\delta\rVert/\lVert x\rVert\) to at most \(5.9\times 10^{-12}\), and the raw residual returns an error below \(1.5\times 10^{-12}\);
  \item a shift of size \(40\lVert x\rVert\) along a kept mode exceeds \(\sqrt{2}\) on every system.
\end{itemize}

Stuck counts at \(2\times 10^{6}\) draws:

\begin{center}
\begin{tabular}{@{}llrr@{}}
\toprule
\(N\) & \(k\) & smallest & energy \\
\midrule
2 & 1 & 0 (0\%) & 0 (0\%) \\
4 & 1 & 8{,}152 (0.41\%) & 381{,}823 (19.1\%) \\
4 & 2 & 8 (0.0004\%) & 74{,}534 (3.73\%) \\
8 & 1 & 23{,}956 (1.20\%) & 738{,}912 (36.9\%) \\
8 & 2 & 138 (0.0069\%) & 338{,}971 (16.9\%) \\
16 & 1 & 56{,}014 (2.80\%) & 1{,}074{,}123 (53.7\%) \\
16 & 2 & 834 (0.042\%) & 672{,}049 (33.6\%) \\
\bottomrule
\end{tabular}
\end{center}

At \(N=2\), removing one mode leaves a single eigenvalue, so \(\kappa_{\mathrm{eff}}=1\).
At \(N=16\) and \(k=1\), the energy rule leaves the condition number essentially unchanged on more than half of the draws.
The rates do not fall when the sample grows from \(10^{6}\) to \(2\times 10^{6}\).

\subsection{Answer}
The large tests verify the classical split of the hybrid method when \(z\) is the exact eigen-piece and the residual solver is exact.
They verify the reduced-mode error identity and the recovery identity of the raw residual in exact arithmetic.
They also verify that the \(\sqrt{2}\) ceiling holds for that orthogonal split and fails for a cancelling explanation.

They do not verify the HHL circuit at these sample sizes.
They do not verify a learned \(z\).
They do not verify a quantum speedup or a wall-clock advantage over a dense solve.
The statevector grid, which does run HHL, stops at \(N=8\) and uses an exact unitary rather than an optimized Hamiltonian simulation.
Removing modes by a full eigendecomposition has already solved \(Ax=b\).

\bibliographystyle{unsrtnat}
\bibliography{refs}

\end{document}
